\section{Alternative Bases and Exact Normalization}
\label{sec:normalized}

CRT exactness is invariant under a change of coordinates, but numerical
geometry is not.  Monomial coordinates $(1,z)$ in a quadratic quotient do not
display the conjugate pair as an orthogonal Euclidean plane.  The basis question
is therefore intrinsic: find a real quotient coordinate whose multiplication
law is the unit complex law.

\subsection{The unit quotient coordinate}

For $0<\theta<\pi$, let
\begin{equation}
 q_\theta=z^2-2\cos\theta\,z+1,
 \qquad
 e_\theta={z-\cos\theta\over\sin\theta}
 \quad\text{in }\K_N[z]/(q_\theta).
 \label{eq:normalized-generator}
\end{equation}

\begin{lemma}[Quadratic unit]
The normalized generator satisfies $e_\theta^2=-1$.
\end{lemma}
\begin{proof}
Writing $c=\cos\theta$ and $s=\sin\theta$, the quotient relation gives
$z^2=2cz-1$.  Therefore
\[
 e_\theta^2={z^2-2cz+c^2\over s^2}
 ={c^2-1\over s^2}=-1.
\]
\end{proof}

\begin{proposition}[Uniqueness among affine generators]
If a real affine class $\eta=\alpha z+\beta$ satisfies
$\eta^2=-1$ in $\R[z]/(q_\theta)$, then
\begin{equation}
 \eta=\pm{z-\cos\theta\over\sin\theta}.
 \label{eq:affine-generator-unique}
\end{equation}
\end{proposition}
\begin{proof}
Reducing $\eta^2$ by $z^2=2cz-1$, the coefficient of $z$ is
$2\alpha(\alpha c+\beta)$.  The case $\alpha=0$ has no real solution to
$\eta^2=-1$, so $\beta=-\alpha c$.  The constant term is then
$-\alpha^2(1-c^2)=-\alpha^2\sin^2\theta$, which equals $-1$ precisely for
$\alpha=\pm1/\sin\theta$.
\end{proof}

\begin{proposition}[Canonical complex chart]
Equip the quotient with the oriented inner product
\begin{equation}
 \langle a+be_\theta,a'+b'e_\theta\rangle_\theta=aa'+bb'.
 \label{eq:quotient-inner-product}
\end{equation}
The map
\begin{equation}
 \phi_\theta:a+be_\theta\longmapsto a+ib
 \label{eq:complex-chart}
\end{equation}
is an isometric real-algebra isomorphism from
$\K_N[z]/(q_\theta)$, after scalar extension to $\R$, onto $\C$.  At the root
$z=e^{-i\theta}$, evaluation is $a+be_\theta\mapsto a-ib$.
\end{proposition}
\begin{proof}
The relation $e_\theta^2=-1$ proves multiplicativity and bijectivity.
The coefficient norm $a^2+b^2$ is the complex modulus squared.  Finally,
$(e^{-i\theta}-\cos\theta)/\sin\theta=-i$.
\end{proof}

For comparison, a raw residue $u+vz$ has normalized coordinates
$(A,B)=(u+v\cos\theta,v\sin\theta)$.  The coordinate matrix
\begin{equation}
 J_\theta=\begin{bmatrix}1&\cos\theta\\0&\sin\theta\end{bmatrix}
 \label{eq:raw-unit-change}
\end{equation}
has singular values $\sqrt{1\pm|\cos\theta|}$ and therefore
\begin{equation}
 \kappa_2(J_\theta)=
 \sqrt{{1+|\cos\theta|\over1-|\cos\theta|}}.
 \label{eq:raw-condition}
\end{equation}
For $0<\theta\le\pi/2$ this is $\cot(\theta/2)$.  The raw monomial chart thus
becomes singular as conjugate roots approach one another, whereas the unit
chart has condition number one by construction.

The choice of $e^{-i\theta}$ fixes the orientation; choosing its conjugate
reverses the sign of the second coordinate.

\subsection{The normalized four-real cell}

The normalized merge of two conjugate-pair states can now be written entirely
over $\R$.  Set
\begin{equation}
 \begin{bmatrix}r\\ i\end{bmatrix}
 =
 \begin{bmatrix}c&-s\\s&c\end{bmatrix}
 \begin{bmatrix}o_a\\o_b\end{bmatrix},
 \quad c^2+s^2=1,
 \label{eq:plane-rotation}
\end{equation}
and define
\begin{equation}
 T_\theta
 \begin{bmatrix}e_a\\e_b\\o_a\\o_b\end{bmatrix}
 =
 \begin{bmatrix}e_a+r\\e_b+i\\e_a-r\\i-e_b\end{bmatrix}.
 \label{eq:normalized-cell}
\end{equation}

\begin{theorem}[Local scaled orthogonality]
For every real $\theta$,
\begin{equation}
 T_\theta^{\trans}T_\theta=2I_4,
 \qquad
 T_\theta^{-1}={1\over2}T_\theta^{\trans}.
 \label{eq:cell-orthogonality}
\end{equation}
The binomial map $H_2(a,b)=(a+b,a-b)$ obeys the same identity.
\end{theorem}
\begin{proof}
The rotation in \eqref{eq:plane-rotation} preserves
$o_a^2+o_b^2$.  Expanding the two add-subtract pairs in
\eqref{eq:normalized-cell} gives
\[
 \|T_\theta(e_a,e_b,o_a,o_b)\|_2^2
 =2(e_a^2+e_b^2+r^2+i^2)=2\|x\|_2^2.
\]
Polarization yields the matrix identity.  The binomial case is the ordinary
Hadamard identity.
\end{proof}

\begin{proposition}[Pair-merge interpretation]
Identify a real pair $(a,b)$ with $a-ib$.  Then
\eqref{eq:normalized-cell}, up to the displayed orientation of the second high
coordinate, is the real-coordinate form of
\begin{equation}
 (E,O)\longmapsto(E+e^{-i\theta}O,\;E-e^{-i\theta}O).
 \label{eq:complex-pair-merge}
\end{equation}
Thus the cell is an exact coordinate expression of a Fourier pair merge.
\end{proposition}
\begin{proof}
Multiplication of $O=o_a-io_b$ by $e^{-i\theta}=c-is$ gives
$(co_a-so_b)-i(so_a+co_b)=r-ii$.  Addition and subtraction from
$E=e_a-ie_b$ give the four entries in
\eqref{eq:normalized-cell}, with the high pair stored in the declared
conjugate orientation.
\end{proof}

\subsection{What the alternative-basis inquiry established}

Two distinct alternative-basis questions sharpened this conclusion.  First,
balanced partitions of conjugate roots were compared by the sparsity of their
CRT reduction projectors.  Circularly adjacent partitions can be dense, while
the covering-compatible partition in \eqref{eq:half-angle-factorization}
retains binomial and trinomial descendants.  This observation motivates the
covering-map family used here; it is not a classification of all sparse CRT
trees.  The June 14--16 radix-four study, conducted after the first DIT
construction, expressed quartic residues in Chebyshev even and odd coordinates.
An isolated quartic node used fewer distinct scalar products, but its state was
not the uniform quadratic two-plane used by neighboring levels.  The normalized basis
\eqref{eq:normalized-generator} is the exact point at which algebraic closure,
unit metric, and a uniform inverse coincide.

\begin{table}[t]
\caption{Coordinates in a terminal real quadratic quotient}
\label{tab:bases}
\centering
\footnotesize
\begin{tabularx}{\columnwidth}{@{}lXX@{}}
\toprule
Basis & Defining relation & Evaluation at $e^{-i\theta}$\\
\midrule
$(1,z)$ & $z^2=2cz-1$ & $a+b(c-is)$\\
$(1,e_\theta)$ & $e_\theta^2=-1$ & $a-ib$\\
\bottomrule
\end{tabularx}
\end{table}

Global orthogonality will follow after the leaf coordinate map has been defined
and identified with Fourier evaluation.  This order avoids assigning a uniform
depth to the binomial spine: different dyadic frequency shells leave that spine
at different levels.
